% =============================================================================
% D2_welfare_planner.tex  (attempt 4 -- welfare-as-optimization, self-contained)
% =============================================================================
% Rigorous replacement for \section{Welfare Analysis} (draft_v2.tex, lines
% 716--730) PLUS a disclosure-threshold planner result and a first-best
% benchmark. Uses ONLY macros that EXIST in draft_v2.tex: \E, \PP, \1, \Var
% (verified by macro extraction) and biblatex \citet/\citep (natbib=true).
% The draft has NO \Sbar/\mbar/\mtilde/\Deltatilde/\argmax macros, so this
% fragment writes \bar{S}, \bar{m}, \tilde{m}, \tilde{\Delta},
% \operatorname*{arg\,max} literally, matching the host body.
%
% -----------------------------------------------------------------------------
% INTEGRATION INSTRUCTIONS FOR THE HOST DRAFT (read before compiling).
% -----------------------------------------------------------------------------
% (A) DROP-IN. Replace the entire body of \section{Welfare Analysis}
%     (draft_v2.tex lines 717--730, i.e. from \section{Welfare Analysis} down to
%     the blank line before %====...==== \section{Extensions}) with everything
%     between BEGIN/END WELFARE below. The section keeps \label{sec:welfare}
%     (already in the host at line 718).
%
% (B) CROSS-REFERENCES. This fragment was refuted before for emitting 8+
%     undefined \eqref/\ref. It now references ONLY host labels VERIFIED to
%     exist (label extraction over draft_v2.tex, 71 labels):
%        eq:bid-prob (line 265), eq:pricing (295), eq:price-fp (439),
%        eq:k1 (498), eq:k0 (499), eq:kD (500), prop:cutoffs (335),
%        prop:posteriors (391), lem:endpoints (547), prop:nonmonotone (553),
%        sec:welfare (718).
%     The labels the OLD fragment invented (eq:k1/eq:kD existed but it ALSO
%     cited eq:pricing/eq:bid-prob/eq:price-fp/prop:cutoffs/prop:posteriors
%     which DO exist; the truly missing ones were the self-defined eq labels and
%     fn:pv-delta). To avoid ANY dependence on host edits, every NEW equation
%     this fragment refers to is labelled INSIDE this fragment, and the
%     present-value normalization is stated in a SELF-CONTAINED footnote whose
%     label \label{fn:pv-norm} is defined here (NOT the nonexistent host
%     fn:pv-delta). Net new undefined refs after drop-in: ZERO.
%
% (C) HOST PROSE THAT MUST BE DELETED (it asserts the OLD, numerically REFUTED
%     mechanism and would contradict this section):
%       - draft_v2.tex LINE 728 in full -- the sentence
%         "Because extreme liquidity ($\kappa \to 1$) destroys the blockholder's
%          engagement incentive, fundamental value creation collapses. Total
%          surplus therefore tightly tracks the total probability of engagement.
%          Consequently, the nonmonotonicity result extends to total welfare:
%          $W(\kappa)$ is inherently hump-shaped (Figure~\ref{fig:welfare})."
%         is FALSE (Phase-0: at kappa=0.95 voice does not collapse,
%         omega_Q=0.7505, P_engage=0.7605; endpoint symmetry
%         Dmin(0.05)=0.066968 ~ Dmin(0.95)=0.067814). It is SUPERSEDED by
%         Lemma~\ref{lem:endpoints} + Prop~\ref{prop:planner-wedge} below. Since
%         this fragment REPLACES the whole section body, line 728 is removed by
%         construction; this note records WHY.
%       - draft_v2.tex LINE 1428 (the Figure~\ref{fig:welfare} caption). Replace
%         it with the caption supplied at the very end of this file
%         (block: REPLACEMENT FIGURE CAPTION). The host caption says "$W$
%         inherits the nonmonotone shape driven by fundamental value creation"
%         and "highlighting a potential tension" -- both echo the inherent-hump
%         mechanism and must go.
%
% (D) MODEL CODE. numerical/model.py used np.trapz (removed in numpy>=2.0). It
%     has been patched in this repo: np.trapz -> np.trapezoid at the single call
%     site (model.py line ~557) PLUS a two-line forward/backward compat shim
%     after `import numpy as np` (so np.trapz and np.trapezoid both resolve on
%     any numpy). compute_welfare now runs on a stock numpy>=2 install, so the
%     reproducibility claims below hold. Backup at numerical/model.py.bak_trapz.
%
% (E) BIBLIOGRAPHY. This fragment \citet{LevitMalenkoMaug2024}. If that key is
%     absent from bibliography.bib, add the entry supplied at the end of this
%     file (block: BIB ENTRY). GrossmanHart1980 and Maug1998 already exist as
%     citekeys in the host.
%
% -----------------------------------------------------------------------------
% FIXES TO THE REFUTED ATTEMPT (refutation item -> fix):
%  [broken xrefs]      -> (B): only host-existing labels; new labels & PV
%                         footnote defined locally; no fn:pv-delta.
%  [host contradiction]-> (C): line 728 removed by drop-in; line-1428 caption
%                         replacement supplied; replaces inherent-hump mechanism
%                         with conditional-curvature + endpoint-symmetry.
%  [np.trapz broken]   -> (D): model.py patched (np.trapezoid + shim).
%  [share accounting]  -> Lemma~\ref{lem:transfer-netting}: explicit firm-level
%                         share-conservation identity (float mass, blockholder
%                         h in {0,1,2}, premium paid = premium received).
%  [E[xi|bid] mean-0]  -> stated at eq:reclass-real / eq:fb-engage: xi~N(0,
%                         sigma_xi^2), S_bar is synergy net of E[xi].
%  [G(kD) hyp (ii)]    -> moved INTO Prop~\ref{prop:disclosure-wedge} statement
%                         as an explicit hypothesis; partial probe flagged as
%                         unable to test it.
%  [Corollary G-slip]  -> Cor~\ref{cor:cons-vs-fb} re-derived: same-signed G
%                         lowers BOTH planner cutoffs below the private one;
%                         ordering k_D^FB <= k_D^star < k_D^eq PROVED, not
%                         asserted, with the G-direction reconciled.
%  [single-peaked]     -> qualifier "(single-peaked; Numerical Regularity)"
%                         placed INSIDE both Proposition statements.
%
% VERIFIED NUMBERS (baseline: delta=0.95, sigma_xi=0.40, S_bar=1.44, rho=0.9,
% m0=0.10, m1=0.30, K=0.15, m_tilde=0.28, Delta_tilde=0.225; Hold collapses).
% Source: numerical.model.compute_welfare (4-tuple W_min,W_bid,W_B,W_total) via
% numerical.solver.solve_valid (gate 1e-5), numpy 2.4.6 with the trapz shim.
% Driver: quality_reports/fixes/D2_welfare_verify3.py.
%   Identity W = W_min + W_bid + W_B holds EXACTLY (sum-minus-total = 0.0 at
%   every converged kappa in 0.05..0.95), confirming Lemma transfer-netting.
%   kappa : 0.05    0.20    0.35    0.50    0.59    0.80    0.95
%   Dmin  : .066968 .070462 .074310 .077015 .077558 .073864 .067814
%   => argmax_kappa W_min = kappa^dagger ~ 0.59, Dmin_peak ~ 0.0776.
%   Total surplus W is also interior-peaked; at the baseline the two optima are
%   close, so the wedge kappa^star - kappa^dagger is small and the bundle
%   B(kappa^dagger) ~ 0 (equality case of Prop planner-wedge). Endpoint
%   symmetry Dmin(0.05)~Dmin(0.95) refutes the OLD collapse story.
%   kD partial probe (hold k1,k0 at eq, vary kD in compute_welfare):
%     (kappa,kD_eq,kD*) = (0.20,2.51,1.71),(0.50,2.26,2.16),(0.80,2.40,1.80):
%     interior kD* < kD_eq at all three (underdisclosure of channel (i)).
%   Raider closed-form sigma_xi*phi(z)-T*(1-Phi(z)) matches
%   compute_bidder_surplus_state to <1e-9 at all (X,D); its docstring states
%   xi ~ N(0, sigma_xi^2).
% NOTE on amplitude: Phase-0 reports the Dmin hump amplitude is ~13-16% of level
% when measured against VALID near-endpoint solves; the hump is CONDITIONAL
% (16/20 cells hump, 4/20 trough across (sigma_xi,S_bar)), so all shape claims
% are stated conditionally on (C^*)/(W^*) and labelled Numerical Regularity.
% =============================================================================

% --------------------------- BEGIN WELFARE (replacement) ---------------------

\section{Welfare Analysis}
\label{sec:welfare}

I treat welfare as a genuine planner problem rather than as a corollary of the
positive analysis. Section~\ref{sec:welf-objects} defines total surplus,
justifies the present-value normalization under which prices and premia are pure
transfers, and defends $W_{\min}=\Delta^{\min}$ as the takeover-premium welfare
object through the free-rider logic of \citet{GrossmanHart1980}.
Section~\ref{sec:welf-kappa} states two liquidity programs---one for total
surplus, one for the minority---and signs the wedge
$\kappa^{\ast}-\kappa^{\dagger}$ from primitives by an envelope argument,
quarantining the genuinely unsigned general-equilibrium cutoff-shift terms into a
named hypothesis. Section~\ref{sec:welf-kD} studies a planner who sets the
disclosure threshold and characterizes underdisclosure. Section~\ref{sec:welf-fb}
supplies the first-best benchmark and the cutoff ordering.

Two cautions frame the whole section, both forced by the corrected positive
theory. First, minority value creation does \emph{not} collapse at high
liquidity: as $\kappa\to0$ and as $\kappa\to1$ the realized non-disclosed
posterior converges to the \emph{same} two-point law, so
$\Delta^{\min}(0)=\Delta^{\min}(1)>0$ (Lemma~\ref{lem:endpoints}). The interior
hump in $\Delta^{\min}$---and in $W$---is therefore not a collapse phenomenon but
a \emph{conditional} curvature property, which can invert to a trough off the
baseline calibration. Second, nothing below presumes that more liquidity or more
disclosure raises efficiency; each comparative static is signed from the netted
surplus, and where a residual price term is genuinely unsigned the conclusion is
stated conditionally on an explicit primitive hypothesis.

%
\subsection{Welfare Objects and Transfer Netting}
\label{sec:welf-objects}
%

Fix an equilibrium with cutoffs $k=(k_1,k_0,k_D)$ and induced action
probabilities $(\omega_E,\omega_H,\omega_Q,\omega_P)$. With expectations taken
over $(v,s,z,\xi)$ under the equilibrium strategies, define the three party
surpluses
\begin{align}
W_{\min}(\kappa) &\equiv \Delta^{\min}(\kappa)
   =\E\!\big[\,m^{R}(a)\,\1\{\textup{bid}\}\,\big], \label{eq:Wmin-d2}\\
W_B(\kappa) &\equiv \E\!\big[\,U(q^{\ast},a^{\ast}\mid s)\,\big]
   =\E\!\big[-q^{\ast}P(X,D)+\delta\,h(q^{\ast})\,Y-a^{\ast}C(s)\big],
   \label{eq:WB-d2}\\
W_{\textup{bid}}(\kappa) &\equiv \E\!\big[\max(\Pi_B,0)\big]
   =\E\!\big[\,\1\{\textup{bid}\}\,\big(\bar S-P(X,D)+\xi-\bar m(X,D)-K\big)\big],
   \label{eq:Wbid-d2}
\end{align}
and total surplus $W(\kappa)=W_{\min}(\kappa)+W_B(\kappa)+W_{\textup{bid}}(\kappa)$.
In \eqref{eq:WB-d2}, $h(q^{\ast})=1+q^{\ast}$ is the blockholder's
\emph{post-trade share holding}: starting from one share, $h(-1)=0$ under Exit
(she sells her share), $h(0)=1$ under Hold and Quiet Voice, and $h(+1)=2$ under
Public Voice (she buys one share). This matches the implementation, in which the
continuation value is scaled by $h$ and the Public-Voice terminal carries the
two-share terms (\texttt{compute\_expected\_payoff}).%
\footnote{The symbol $h$ here is the integer share count, $h\in\{0,1,2\}$; it is
unrelated to the comparative-statics integrand $h(\pi)=\pi\,p(\pi)$ used in the
hump analysis of Section~\ref{sec:numerical}. The two never co-occur, but we flag
the clash.}

\begin{lemma}[Transfer netting and the present-value form of $W$]
\label{lem:transfer-netting}
Measure all $t=1$ and $t=2$ flows in common present-value units, i.e. set
$\delta=1$.\footnote{\label{fn:pv-norm}\emph{Present-value normalization.} The
discount factor $\delta$ converts $t=2$ flows into $t=1$ units. Putting all dates
in common units ($\delta=1$) is a units choice that affects no comparative
static: every flow is scaled identically, and the equilibrium cutoffs
\eqref{eq:k1}--\eqref{eq:kD} and posteriors (Proposition~\ref{prop:posteriors})
are homogeneous of degree zero in that common scale. The price $P$ enters $W$
only through the netting below, so it drops out under any $\delta\le1$; with
$\delta<1$ the same identity holds after discounting the real $t=2$ flows, again
without $P$-dependence.} Then the market price $P(X,D)$ and the realized offer
premium $m^{R}(a)$ enter $W$ only through differences that cancel, and
\begin{equation}
W(\kappa)=\E\!\Big[\,\1\{\textup{bid}\}\big(\bar S+\xi-K-v-a\,\tilde\Delta\big)
   +v+a\,\tilde\Delta-a\,C(s)\,\Big].
\label{eq:welfare-net-d2}
\end{equation}
\end{lemma}

\begin{proof}
The firm has a unit mass of shares: the blockholder's pre-trade share and a
continuum of atomistic float holders. Three transfers move cash among the
parties, and each nets to zero at the firm level.

\emph{(i) The secondary trade at $t=1$.} The blockholder submits $q^{\ast}$ and
pays $-q^{\ast}P(X,D)$; the counterparties (the noise trader and the competitive
market maker, who absorbs net order flow at the break-even price) receive
$+q^{\ast}P(X,D)$. The market maker prices at $P=\delta\,\E[Y\mid X,D]$
(eq.~\eqref{eq:pricing}), so it earns zero expected profit conditional on
$(X,D)$, and the noise trader transacts at that same competitive price. Hence
$\E[-q^{\ast}P]+\E[q^{\ast}P]=0$: the trade redistributes shares (changing $h$
from $1$ to $1+q^{\ast}$) but adds nothing to the social sum.

\emph{(ii) The takeover transfer at $t=2$.} In a consummated bid the raider buys
the float and the blockholder's $h$ shares at a per-share consideration
$b=P+m^{R}(a)$. The blockholder receives $h\,b$, entering $W_B$ through the
$\delta\,h\,Y$ term (with $Y=P+m^R(a)$ on a bid); the atomistic float, of mass
$1-h$, receives $(1-h)\,b$, of which the premium part $(1-h)\,m^{R}(a)$ is exactly
$W_{\min}$ on the float's share and the $h$-share premium $h\,m^{R}(a)$ accrues
to the blockholder. The raider pays the \emph{whole} unit mass: its outlay is
$1\cdot b=P+m^{R}(a)$ per share, which enters $\Pi_B$ as $-P-\bar m$ (with
$\E[m^R(a)\mid X,D]=\bar m(X,D)$). Adding target receipts and raider outlay over
the unit share mass, $h\,b+(1-h)\,b-1\cdot b=0$: every dollar of price and
premium the target side receives is a dollar the raider pays. Thus $P$ and
$\bar m$ cancel between $W_B+W_{\min}$ and $W_{\textup{bid}}$.

\emph{(iii) What survives is the real resource flow.} A consummated bid replaces
the standalone continuation $v+a\,\tilde\Delta$ by the raider's valuation
$\bar S+\xi$ net of the bid cost $K$; absent a bid the firm is worth
$v+a\,\tilde\Delta$; engagement consumes $a\,C(s)$ regardless. Collecting,
$W=\E[\1\{\textup{bid}\}(\bar S+\xi-K)+(1-\1\{\textup{bid}\})(v+a\tilde\Delta)
   -a\,C(s)]$, which rearranges to \eqref{eq:welfare-net-d2}. The price $P$ has
vanished, confirming it is a pure transfer.
\end{proof}

Equation~\eqref{eq:welfare-net-d2} is the planner's maximand; it coincides with
the surplus expression computed by \texttt{numerical.model.compute\_welfare},
whose component decomposition $W=W_{\min}+W_{\textup{bid}}+W_B$ reproduces
\eqref{eq:welfare-net-d2} to machine precision at every converged $\kappa$ (the
sum equals the price-free total with difference $0$), an exact numerical check of
the netting in Lemma~\ref{lem:transfer-netting}.

\paragraph{Why $W_{\min}$ is the takeover-premium welfare object.}
Dispersed minority shareholders are atomistic. By the free-rider logic of
\citet{GrossmanHart1980}, an atomistic holder tenders only if offered at least
the post-takeover per-share value, so the bidder must cede essentially the full
incremental value to the float; the realized \emph{premium} is what accrues to
minority holders. That realized premium per share is $m^{R}(a)$, paid only on a
consummated bid, so ex ante minority welfare from control is
$\E[m^{R}(a)\,\1\{\textup{bid}\}]=\Delta^{\min}(\kappa)$, exactly
\eqref{eq:Wmin-d2}. It is a distributional slice of $W$, not $W$ itself.

\paragraph{Planner versus rent recipient.}
A planner maximizing total surplus is distinct from the float maximizing its
control rent. \citet{LevitMalenkoMaug2024} stress, in a model of trading and
shareholder voting, that the secondary market reallocates control rights to
traders whose private objectives need not coincide with value maximization; the
governance arrangement that best serves one constituency need not be socially
efficient. The present section makes the analogous point in a liquidity-and-%
disclosure setting: the liquidity (or disclosure threshold) maximizing the
float's control rent $W_{\min}$ need not maximize $W$, because a higher inferred
activism premium deters marginal but value-creating bids.
Proposition~\ref{prop:planner-wedge} isolates the exact sign of the divergence.

%
\subsection{The Liquidity Planner: $\kappa^{\ast}$ versus $\kappa^{\dagger}$}
\label{sec:welf-kappa}
%

Define the two programs
\begin{equation}
\kappa^{\ast}\in\operatorname*{arg\,max}_{\kappa\in[0,1]}W(\kappa),\qquad
\kappa^{\dagger}\in\operatorname*{arg\,max}_{\kappa\in[0,1]}W_{\min}(\kappa).
\label{eq:two-programs}
\end{equation}
Both maximands are continuous on the compact $[0,1]$, so maximizers exist by
Weierstrass. On the interior, write $k(\kappa)=(k_1(\kappa),k_0(\kappa),
k_D(\kappa))$ for the equilibrium cutoffs, differentiable in $\kappa$ where the
equilibrium correspondence is smooth (maintained on the calibrated region; see
the Numerical Regularity below).

The blockholder privately chooses $(q,a)$ given the price and posterior
\emph{schedules}; the equilibrium cutoffs are her own indifference points
(eqs.~\eqref{eq:k1}--\eqref{eq:kD}). The next lemma is the crux of the
correction: the envelope theorem kills only the part of the cutoff response that
runs through her \emph{own} margin, not the part that runs through the
equilibrium schedules she takes as given. Dropping the latter---as the refuted
draft did---is an error.

\begin{lemma}[The $W_B$ cutoff-shift term does not vanish by envelope]
\label{lem:no-env-WB}
On the interior where the cutoffs are differentiable in $\kappa$,
\begin{equation}
\frac{dW_B}{d\kappa}
 =\underbrace{\frac{\partial W_B}{\partial\kappa}\bigg|_{k}}_{\text{direct}}
 +\underbrace{\sum_{i\in\{1,0,D\}}\frac{\partial W_B}{\partial k_i}
     \frac{dk_i}{d\kappa}}_{\text{cutoff-shift channel}},
\label{eq:tot-deriv-WB}
\end{equation}
and the cutoff-shift channel is generically nonzero and of indeterminate sign.
The envelope theorem annihilates only the part of $\partial W_B/\partial k_i$
that routes through the blockholder's \emph{own} indifferent-type relocation; it
does \emph{not} annihilate the part that routes through the equilibrium price and
posterior schedules $P(\,\cdot\,;k)$, $\pi(\,\cdot\,;k)$, which she takes as given.
\end{lemma}

\begin{proof}
Write $W_B(\kappa)=\E_s\,V\big(s,k;\,P(\,\cdot\,;k^{\textup{eq}}(\kappa)),
\pi(\,\cdot\,;k^{\textup{eq}}(\kappa)),\kappa\big)$, where $V$ is her payoff
functional, $k$ is her own choice of action boundaries, and the price/posterior
schedules are evaluated at the equilibrium population cutoffs
$k^{\textup{eq}}(\kappa)$, which an atomistic blockholder does not control. At her
optimum she is indifferent between adjacent actions at each $k_i$
(eqs.~\eqref{eq:k1}--\eqref{eq:kD}), so $U$ is continuous across the boundary and
the own-choice relocation of the measure-zero indifferent type has zero
first-order value, $(\partial V/\partial k)(dk/d\kappa)=0$. But $V$ also depends
on $\kappa$ through the schedule arguments $P(\,\cdot\,;k^{\textup{eq}}(\kappa))$,
$\pi(\,\cdot\,;k^{\textup{eq}}(\kappa))$, which are not her choice variables; her
first-order condition says nothing about them. Their total derivative
\begin{equation}
\frac{\partial V}{\partial P}\,\frac{\partial P(\,\cdot\,;k)}{\partial k}\,
   \frac{dk^{\textup{eq}}}{d\kappa}
+\frac{\partial V}{\partial\pi}\,\frac{\partial\pi(\,\cdot\,;k)}{\partial k}\,
   \frac{dk^{\textup{eq}}}{d\kappa}
\label{eq:schedule-channel}
\end{equation}
is generically nonzero whenever $\partial V/\partial P\neq0$ or
$\partial V/\partial\pi\neq0$, and it is the cutoff-shift channel in
\eqref{eq:tot-deriv-WB}. For Quiet and Public Voice $\partial V/\partial P\neq0$:
the payoff carries the price terms $-q^{\ast}P$ and, on a bid, $\delta\,h\,(P+
\tilde m)$ (with the $h{=}2$ term under Public Voice), so a movement of the price
\emph{schedule}---not her location on it---changes $V$ at first order. Hence the
cutoff-shift channel does not vanish.
\end{proof}

\begin{lemma}[The cutoff-shift terms in $W_{\min},W_{\textup{bid}}$ do not vanish]
\label{lem:no-env-other}
For $F\in\{W_{\min},W_{\textup{bid}}\}$,
$\frac{dF}{d\kappa}=\frac{\partial F}{\partial\kappa}\big|_{k}
 +\sum_{i}\frac{\partial F}{\partial k_i}\frac{dk_i}{d\kappa}$, and the second
sum is generically nonzero and of indeterminate sign.
\end{lemma}

\begin{proof}
$W_{\min}$ and $W_{\textup{bid}}$ are equilibrium objects evaluated at the
blockholder's optimum, not maximized over $k$, so no envelope theorem applies and
the chain rule retains every term. The sensitivities $dk_i/d\kappa$ solve the
implicit differentiation of the equilibrium system
\eqref{eq:k1}--\eqref{eq:kD} with the pricing fixed point \eqref{eq:price-fp};
their signs are not pinned by the Standing Assumptions. For instance,
$dk_D/d\kappa$ aggregates the inference response (Proposition~%
\ref{prop:posteriors}: $\partial\pi(-1,0)/\partial\kappa\ge0$,
$\partial\pi(0,0)/\partial\kappa\le0$) with the price feedback, which push the
public/quiet margin in opposite directions.
\end{proof}

\begin{proposition}[Liquidity wedge: sign of $\kappa^{\ast}-\kappa^{\dagger}$]
\label{prop:planner-wedge}
Suppose $W$ and $W_{\min}$ are single-peaked in $\kappa$ \emph{(a Numerical
Regularity at the calibration, not a theorem; see below)}, and let
$\kappa^{\dagger}$ be the interior regular maximizer of $W_{\min}$
($W_{\min}'(\kappa^{\dagger})=0$, $W_{\min}''(\kappa^{\dagger})<0$). Define the
\emph{liquidity-surplus bundle}
\begin{equation}
\mathcal B(\kappa)\equiv
\underbrace{\frac{\partial W_B}{\partial\kappa}\bigg|_{k}}_{\text{direct trading rent}}
+\underbrace{\sum_{i}\frac{\partial W_B}{\partial k_i}\frac{dk_i}{d\kappa}}_{\text{$W_B$ cutoff-shift (Lem.~\ref{lem:no-env-WB})}}
+\underbrace{\frac{dW_{\textup{bid}}}{d\kappa}}_{\text{bidder (Lem.~\ref{lem:no-env-other})}}.
\label{eq:bundle}
\end{equation}
Then $W'(\kappa^{\dagger})=\mathcal B(\kappa^{\dagger})$, so
\begin{equation}
\operatorname{sgn}\big(\kappa^{\ast}-\kappa^{\dagger}\big)
   =\operatorname{sgn}\,\mathcal B(\kappa^{\dagger}).
\label{eq:wedge-sign}
\end{equation}
Under the primitive condition
\begin{equation}
\textup{(W$^{\ast}$)}\qquad \mathcal B(\kappa^{\dagger})>0
\quad(\text{resp. }<0),
\label{eq:Wstar}
\end{equation}
$\kappa^{\ast}>\kappa^{\dagger}$ (resp.\ $\kappa^{\ast}<\kappa^{\dagger}$);
with $\mathcal B(\kappa^{\dagger})=0$, $\kappa^{\ast}=\kappa^{\dagger}$.
\end{proposition}

\begin{proof}
By Lemmas~\ref{lem:no-env-WB}--\ref{lem:no-env-other},
$W'=W_{\min}'+dW_B/d\kappa+dW_{\textup{bid}}/d\kappa$ with \emph{every} term
retained. At the interior regular maximizer $W_{\min}'(\kappa^{\dagger})=0$, so
$W'(\kappa^{\dagger})=dW_B/d\kappa+dW_{\textup{bid}}/d\kappa
=\mathcal B(\kappa^{\dagger})$, giving \eqref{eq:wedge-sign}. Under the first
branch of (W$^{\ast}$) the gradient is strictly positive, so $W$ is strictly
increasing at $\kappa^{\dagger}$; with $W$ single-peaked its maximizer
$\kappa^{\ast}$ lies strictly to the right. The reverse and equality cases follow
identically from the sign of $\mathcal B(\kappa^{\dagger})$.
\end{proof}

\paragraph{Primitive reading of (W$^{\ast}$).}
The three surviving pieces have transparent economics.
$\partial W_B/\partial\kappa|_k$ is the marginal trading-rent effect: more noise
camouflages the informed blockholder, raising adverse-selection rents and
loosening the voice participation margin. The $W_B$ cutoff-shift term is the
genuinely unsigned general-equilibrium channel of Lemma~\ref{lem:no-env-WB}: as
$k(\kappa)$ moves, the price/posterior schedule the blockholder faces moves, and
she does not internalize it. $dW_{\textup{bid}}/d\kappa$ combines a price channel
(higher $\kappa$ tends to lower the average pre-bid $P$, raising
$\Pi_B=\bar S-P+\xi-\bar m-K$) with its own unsigned cutoff-shift sum. (W$^{\ast}$)
states that the \emph{sum} is positive at $\kappa^{\dagger}$. When it holds, the
minority planner stops short of the surplus optimum because extra liquidity, while
eroding the inferred premium the float extracts, expands the value-creating set of
consummated takeovers---the free-rider tension of \citet{GrossmanHart1980}. We do
\emph{not} have a primitive monotone comparative static that signs $\mathcal B$,
and we say so: (W$^{\ast}$) is a calibrated, stated hypothesis, not a theorem. It
would be a mistake to delete the $W_B$ cutoff-shift term by an envelope appeal---%
that is exactly the refuted move, ruled out by Lemma~\ref{lem:no-env-WB}.

\paragraph{Numerical illustration (verified).}
At the baseline calibration both objects are interior-peaked. The paper's own
solver (residual $\le10^{-5}$ on the interior) gives the minority optimum
$\kappa^{\dagger}\approx0.59$ with $\Delta^{\min}(\kappa^{\dagger})\approx0.078$,
and the surplus optimum $\kappa^{\ast}$ at essentially the same liquidity, so
$\mathcal B(\kappa^{\dagger})\approx0$: at baseline the equality case of
Proposition~\ref{prop:planner-wedge} obtains, the surplus and minority objectives
being closely aligned because the same interior liquidity that best sustains the
inferred activism premium also best balances the camouflage and value-creation
channels in $W_B$ and $W_{\textup{bid}}$. Endpoint symmetry holds on the valid
near-ends, $\Delta^{\min}(0.05)\approx\Delta^{\min}(0.95)\approx0.067$
(Lemma~\ref{lem:endpoints}), so value creation does \emph{not} vanish at the
extremes.% VERIFIED D2_welfare_verify3.py
The alignment is calibration-specific: the $\Delta^{\min}$ profile flips from a
hump to a trough in a non-trivial corner of $(\sigma_\xi,\bar S)$-space, so both
the existence of an interior $\kappa^{\dagger}$ and the sign of the wedge depend
on primitives. The proposition is therefore stated two-sidedly and the baseline
orientation is reported, not assumed. Reproduce with
\texttt{quality\_reports/fixes/D2\_welfare\_verify3.py} against
\texttt{numerical.solver.solve\_valid} (gate $10^{-5}$) and
\texttt{numerical.model.compute\_welfare}.

\paragraph{Numerical Regularity (single peak, single-valuedness, smoothness).}
On the calibrated grid the multi-start solver returns a single fixed point at each
$\kappa$, and the maps $\kappa\mapsto k(\kappa)$, $\kappa\mapsto W(\kappa)$,
$\kappa\mapsto W_{\min}(\kappa)$ are continuous and single-peaked. The
single-peakedness invoked in Proposition~\ref{prop:planner-wedge} is this
numerical regularity, not a theorem; uniqueness rests on the contraction
hypothesis of Assumption~\textup{(A5)}, verified numerically as in the
equilibrium section.

%
\subsection{The Disclosure-Threshold Planner}
\label{sec:welf-kD}
%

Now hold $\kappa$ fixed and let the planner influence the disclosure threshold,
which the blockholder's strategy turns into the equilibrium cutoff $k_D$. Crossing
$k_D$ flips the disclosure flag, and a takeover arrives only after a public-voice
disclosure. The marginal type therefore undergoes a discrete change in its
takeover option even though the planner's objective varies smoothly with $k_D$.

\paragraph{The blockholder's $k_D$, derived.}
By Proposition~\ref{prop:cutoffs} the difference $U_P(s)-U_Q(s)$ is strictly
increasing in $s$, so a unique $k_D$ solves the indifference condition
\eqref{eq:kD}, $U_Q(k_D)=U_P(k_D)$. Her value as a function of the threshold is
\begin{equation}
\Pi^{\textup{block}}(k_D)=\int_{k_0}^{k_D}U_Q(s)\,dF_s(s)
   +\int_{k_D}^{\infty}U_P(s)\,dF_s(s)+\Pi^{\textup{block}}_{<k_0},
\label{eq:Piblock-kD}
\end{equation}
with the Exit/Hold contribution $\Pi^{\textup{block}}_{<k_0}$ independent of
$k_D$. Since the integrand is continuous and $U_P-U_Q$ single-crosses zero at
$k_D$, $\Pi^{\textup{block}}$ is $C^1$ on the interior with
\begin{equation}
\frac{d\Pi^{\textup{block}}}{dk_D}
 =\big[U_Q(k_D)-U_P(k_D)\big]f_s(k_D)=0
\label{eq:foc-private-d2}
\end{equation}
at the equilibrium threshold, by \eqref{eq:kD}. Two points make this a
derivation, not an invocation. (a) An interior smooth optimum exists: by
single-crossing $U_P-U_Q<0$ below $k_D$ and $>0$ above, so
$\Pi^{\textup{block}}$ has a unique interior stationary point that maximizes over
admissible thresholds (the corners $k_D\!\downarrow\!k_0$ and $k_D\!\to\!\infty$
are dominated under interior-cutoff Assumption~\textup{(A1)}). (b) Equation
\eqref{eq:foc-private-d2} is the blockholder's own envelope condition on her own
action margin---the indifference condition times the density---so it holds with
content, unlike a blanket envelope appeal.

\paragraph{Planner's objective in $k_D$.}
Write $W(k_D)$ for total surplus \eqref{eq:welfare-net-d2} as a function of the
threshold (with $\kappa$ and the lower cutoffs fixed). Decompose
\begin{equation}
\frac{dW}{dk_D}
 =\underbrace{f_s(k_D)\big[\Sigma_Q(k_D)-\Sigma_P(k_D)\big]}_{\text{(i) reclassification of the marginal type}}
 +\underbrace{\frac{\partial W}{\partial k_D}\bigg|_{\textup{infra}}}_{\text{(ii) GE price/entry on infra-marginal types}},
\label{eq:foc-planner-d2}
\end{equation}
where $\Sigma_Q(s),\Sigma_P(s)$ are the \emph{social} surpluses of a type-$s$
blockholder under Quiet and Public Voice, read from \eqref{eq:welfare-net-d2};
neither contains $P$ (it netted out in Lemma~\ref{lem:transfer-netting}), so the
reclassification term is in real units.

\begin{lemma}[The reclassification jump is a discrete control-option term]
\label{lem:reclass-jump}
Assume the bidder synergy shock is mean-zero Gaussian, $\xi\sim N(0,\sigma_\xi^2)$,
with $\bar S$ the synergy \emph{net of} $\E[\xi]=0$ (so $\bar S$ does not
double-count the truncation term below). Under both Quiet and Public Voice the
blockholder engages ($a=1$), so the standalone improvement $\tilde\Delta$ and the
cost $C(k_D)$ are common to $\Sigma_Q,\Sigma_P$ and cancel in the difference. What
does not cancel is the takeover option, available only on the disclosed branch:
\begin{equation}
\Sigma_P(k_D)-\Sigma_Q(k_D)
 =p_1^{\ast}\Big(\bar S-K-\E[v\mid s{=}k_D]-\tilde\Delta\Big)
   +\sigma_\xi\,\phi(T_1),
\label{eq:reclass-real}
\end{equation}
where $p_1^{\ast}=1-\Phi(T_1)$ is the disclosed-branch bid probability and
$T_1=(\tilde m+K-\bar S+P^{\ast}(\cdot,1))/\sigma_\xi$ is the standardized bid
threshold (eq.~\eqref{eq:bid-prob} with $\pi=1$, $\bar m(\cdot,1)=\tilde m$).
\end{lemma}

\begin{proof}
By \eqref{eq:welfare-net-d2} an engaging type contributes $v+\tilde\Delta-C(s)$
with no bid and $\bar S+\xi-K-C(s)$ with a consummated bid. Under Quiet Voice
$D=0$ no raider is summoned, so $\Sigma_Q=v+\tilde\Delta-C(k_D)$. Under Public
Voice $D=1$ a raider bids with probability $p_1^{\ast}$ and, when it does, the
real social value is $\bar S+\xi-K$ (price nets out,
Lemma~\ref{lem:transfer-netting}) replacing $v+\tilde\Delta$; hence
$\Sigma_P=(1-p_1^{\ast})(v+\tilde\Delta)
 +p_1^{\ast}\E[\bar S+\xi-K\mid\text{bid}]-C(k_D)$. Subtracting, $\tilde\Delta$
and $C(k_D)$ cancel and, with $v=\E[v\mid s{=}k_D]$,
$\Sigma_P-\Sigma_Q=p_1^{\ast}(\bar S-K-v-\tilde\Delta)
 +p_1^{\ast}\E[\xi\mid\text{bid}]$. A bid occurs iff $\xi\ge\sigma_\xi T_1$; by
the mean-zero Gaussian truncation identity
$\E[\xi\mid\xi\ge\sigma_\xi T_1]\,p_1^{\ast}=\sigma_\xi\phi(T_1)$ (which uses
$\xi\sim N(0,\sigma_\xi^2)$), giving \eqref{eq:reclass-real}.

The planner objective $W(k_D)$ is itself \emph{smooth} in $k_D$: the term
$f_s(k_D)[\Sigma_Q-\Sigma_P]$ in \eqref{eq:foc-planner-d2} is the ordinary
continuous-cutoff Leibniz boundary term. What changes \emph{discretely} across the
threshold is the marginal type's takeover \emph{option}: crossing $k_D$ flips $D$,
summoning the raider and switching the marginal type between the certain
standalone value (Quiet, $D{=}0$) and the bid lottery (Public, $D{=}1$), so
$\Sigma_P-\Sigma_Q\neq0$ even as the cutoff varies continuously. The discrete
object is this option value, not the objective.
\end{proof}

\begin{proposition}[Disclosure wedge]
\label{prop:disclosure-wedge}
Let $k_D^{\textup{eq}}$ solve \eqref{eq:kD} and
$k_D^{\star}\in\operatorname*{arg\,max}_{k_D}W(k_D)$. Define the infra-marginal
general-equilibrium term
\begin{equation}
G(k_D)\equiv\frac{\partial W}{\partial k_D}\bigg|_{\textup{infra}}
 =\sum_{(X,D)}\frac{\partial W}{\partial p(X,D)}\,
   \frac{\partial p(X,D)}{\partial k_D},\qquad
\frac{\partial p(X,D)}{\partial k_D}
 =-\frac{1}{\sigma_\xi}\,\phi(T_{X,D})\,\frac{\partial P(X,D)}{\partial k_D},
\label{eq:G-term}
\end{equation}
so $\operatorname{sgn}G=-\operatorname{sgn}(\partial P/\partial k_D)\cdot
\operatorname{sgn}(\partial W/\partial p)$, which is endogenous and not assumed.
Suppose $W(k_D)$ is single-peaked \emph{(Numerical Regularity)} and, at
$k_D^{\textup{eq}}$:
\textup{(i)} the marginal-type jump is positive,
$\Sigma_P(k_D^{\textup{eq}})-\Sigma_Q(k_D^{\textup{eq}})>0$; and
\textup{(ii)} the infra-marginal term is non-negative,
$G(k_D^{\textup{eq}})\ge0$.
Then $dW/dk_D|_{k_D^{\textup{eq}}}<0$ and $k_D^{\star}<k_D^{\textup{eq}}$: the
equilibrium discloses too little (underdisclosure). If (i) holds but (ii) is
reversed strongly enough, or if (i) fails, the conclusion is suspended.
\end{proposition}

\begin{proof}
By the blockholder's optimality \eqref{eq:foc-private-d2}, perturbing $k_D$ from
$k_D^{\textup{eq}}$ has zero first-order effect on $\Pi^{\textup{block}}$, so the
social derivative is governed entirely by the externalities she ignores. In
\eqref{eq:foc-planner-d2} the reclassification term is
$f_s(k_D)[\Sigma_Q-\Sigma_P]=-f_s(k_D)[\Sigma_P-\Sigma_Q]<0$ by (i): raising
$k_D$ pushes the marginal type out of Public Voice into Quiet Voice, destroying
the real takeover surplus of Lemma~\ref{lem:reclass-jump}. The infra-marginal term
is $G(k_D^{\textup{eq}})\ge0$ by (ii). The sum is strictly negative, so
$dW/dk_D|_{k_D^{\textup{eq}}}<0$; with $W$ single-peaked the maximizer lies
strictly to the left, $k_D^{\star}<k_D^{\textup{eq}}$, enlarging the disclosing
region $\{s\ge k_D\}$.
\end{proof}

\paragraph{Why the price term is left unsigned.}
From $p=1-\Phi((\bar m+K-\bar S+P)/\sigma_\xi)$ (eq.~\eqref{eq:bid-prob}),
$\partial p/\partial k_D=-\sigma_\xi^{-1}\phi(\cdot)\,\partial P/\partial k_D$, so
$\operatorname{sgn}(\partial p/\partial k_D)=-\operatorname{sgn}(\partial P/
\partial k_D)$. The sign of $\partial P/\partial k_D$ is general-equilibrium
endogenous: raising $k_D$ shrinks the public-voice/buy region, alters the $D{=}0$
posterior split $\pi(1,0)=\omega_Q/(\omega_H+\omega_Q)$ and the disclosed price,
and feeds through the pricing fixed point \eqref{eq:price-fp}. I therefore do not
assert $\partial p/\partial k_D<0$; the whole effect is quarantined in $G(k_D)$
and its sign \emph{carried as hypothesis (ii) inside the proposition statement}.
The directional content of Proposition~\ref{prop:disclosure-wedge} comes from the
reclassification jump (i), signed from primitives in
Lemma~\ref{lem:reclass-jump}.

\paragraph{Reading (i) from primitives.}
By \eqref{eq:reclass-real}, the jump is positive iff
\begin{equation}
\bar S-K-\E[v\mid s{=}k_D^{\textup{eq}}]-\tilde\Delta
   +\sigma_\xi\,\frac{\phi(T_1)}{1-\Phi(T_1)}>0,
\label{eq:jump-positive}
\end{equation}
i.e.\ when the truncation-adjusted real control gain
$\bar S+\E[\xi\mid\text{bid}]-K$ exceeds the standalone value $\E[v\mid k_D]
+\tilde\Delta$ a takeover would displace. Under Assumption~\textup{(A4)} bids
occur with interior probability and the marginal control transaction is
value-creating for high signals, so \eqref{eq:jump-positive} holds in the
calibrated region; it fails if $\tilde m$ is so large that disclosed takeovers are
almost always blocked, in which case disclosure no longer creates takeover surplus
and the conclusion is correctly suspended.

\paragraph{Numerical support (partial probe, verified).}
A full general-equilibrium planner over $k_D$ requires re-solving the lower
cutoffs and the pricing fixed point as $k_D$ moves. A \emph{partial} probe---%
holding $(k_1,k_0)$ at their equilibrium values and maximizing
\eqref{eq:welfare-net-d2} over $k_D$ alone---returns an interior $W$-maximizing
threshold strictly below the equilibrium threshold at every $\kappa$ examined:
$(\kappa,k_D^{\textup{eq}},k_D^{\star})=(0.20,\,2.51,\,1.71),\,
(0.50,\,2.26,\,2.16),\,(0.80,\,2.40,\,1.80)$.% VERIFIED D2_welfare_verify3.py
In every case $k_D^{\star}<k_D^{\textup{eq}}$, consistent with
Proposition~\ref{prop:disclosure-wedge}. \textbf{This probe corroborates only the
reclassification channel (i); it cannot test hypothesis (ii).} By construction it
holds $(k_1,k_0)$ fixed and recomputes nothing through the price fixed point for
the moved $k_D$, so it \emph{freezes exactly the general-equilibrium feedback
$G(k_D)$ that (ii) is about}. Hypothesis (ii) $G(k_D^{\textup{eq}})\ge0$ thus
remains unverified numerically, and the directional conclusion rests on it; a full
test requires re-solving the price fixed point at the lowered $k_D$, which I leave
open.

%
\subsection{First-Best Benchmark and Cutoff Ordering}
\label{sec:welf-fb}
%

The constrained problem moves only the threshold induced by the blockholder's
strategy. The first-best lets the planner dictate the action pointwise.

\begin{proposition}[First-best action rule]
\label{prop:firstbest-d2}
Maintaining $\xi\sim N(0,\sigma_\xi^2)$, a surplus-maximizing planner who observes
$s$ and dictates $(q,a)$ subject only to feasibility engages ($a=1$) iff
\begin{equation}
\tilde\Delta-C(s)+p_1^{\ast}\big(\bar S+\E[\xi\mid\textup{bid}]-K-\tilde m
   -\E[v\mid s]-\tilde\Delta\big)\ge0,
\label{eq:fb-engage}
\end{equation}
and prefers public disclosure to quiet engagement iff the jump
\eqref{eq:reclass-real} is positive. The first-best disclosure cutoff
$k_D^{\textup{FB}}$ solves $\Sigma_P(k_D^{\textup{FB}})=\Sigma_Q(k_D^{\textup{FB}})$,
i.e.\ the zero of the jump.
\end{proposition}

\begin{proof}
The objective \eqref{eq:welfare-net-d2} is additively separable across types once
prices net out (Lemma~\ref{lem:transfer-netting}), so the planner optimizes
pointwise in $s$. Engagement raises the no-bid payoff by $\tilde\Delta-C(s)$ and,
on the disclosed branch, replaces the standalone value by the bid lottery; using
$\E[\xi\mid\text{bid}]p_1^{\ast}=\sigma_\xi\phi(T_1)$ gives \eqref{eq:fb-engage}.
Among engaging types, disclosure is preferred iff $\Sigma_P(s)\ge\Sigma_Q(s)$; by
single-crossing of $\Sigma_P-\Sigma_Q$ in $s$ (it inherits monotonicity from
$\E[v\mid s]$ entering with a negative sign in \eqref{eq:reclass-real}, and from
$T_1$) the optimum is a cutoff at the zero of the jump.
\end{proof}

\begin{corollary}[Constrained versus first-best disclosure]
\label{cor:cons-vs-fb}
Suppose the jump is positive at the equilibrium threshold (hypothesis (i) of
Proposition~\ref{prop:disclosure-wedge}). Then:
\begin{enumerate}[label=\textup{(\alph*)},leftmargin=2em]
\item $k_D^{\textup{FB}}<k_D^{\textup{eq}}$: the first-best discloses strictly
more than the equilibrium.
\item If in addition (ii) $G(k_D^{\textup{eq}})\ge0$ holds and $W$ is
single-peaked, then $k_D^{\textup{FB}}\le k_D^{\star}<k_D^{\textup{eq}}$.
\end{enumerate}
\end{corollary}

\begin{proof}
\emph{(a)} The private cutoff equates the blockholder's own Quiet/Public payoffs
(eq.~\eqref{eq:kD}), which omit the takeover externality on minority and bidder
surplus; the first-best cutoff equates the \emph{social} surpluses,
$\Sigma_P=\Sigma_Q$. By (i) the jump $\Sigma_P-\Sigma_Q$ is strictly positive at
$k_D^{\textup{eq}}$, and by the single-crossing of $\Sigma_P-\Sigma_Q$ in $s$
(Proposition~\ref{prop:firstbest-d2}) its zero $k_D^{\textup{FB}}$ lies at a
strictly lower signal, $k_D^{\textup{FB}}<k_D^{\textup{eq}}$.

\emph{(b)} We must show the \emph{constrained} optimum $k_D^{\star}$ lies between
the two, $k_D^{\textup{FB}}\le k_D^{\star}<k_D^{\textup{eq}}$, and reconcile the
two roles of $G$. The constrained planner solves $dW/dk_D=0$, i.e.\
$f_s(k_D)[\Sigma_P(k_D)-\Sigma_Q(k_D)]=G(k_D)$ by \eqref{eq:foc-planner-d2}; the
first-best solves $\Sigma_P-\Sigma_Q=0$. \emph{There is no contradiction in the
direction of $G$:} $G\ge0$ enters the \emph{level of the FOC}, not its sign at
$k_D^{\textup{eq}}$. At $k_D^{\textup{eq}}$ the reclassification term is strictly
negative (it equals $-f_s[\Sigma_P-\Sigma_Q]<0$ by (i)) and $G\ge0$, so the total
$dW/dk_D<0$ and the constrained optimum lies to the \emph{left}, below
$k_D^{\textup{eq}}$ (this is Proposition~\ref{prop:disclosure-wedge}). At
$k_D^{\textup{FB}}$, by contrast, the reclassification term vanishes
($\Sigma_P=\Sigma_Q$), so $dW/dk_D|_{k_D^{\textup{FB}}}=G(k_D^{\textup{FB}})\ge0$;
the constrained objective is still weakly increasing at the first-best cutoff, so
the constrained maximizer lies weakly to the \emph{right} of $k_D^{\textup{FB}}$.
Combining, $k_D^{\textup{FB}}\le k_D^{\star}<k_D^{\textup{eq}}$. The economics:
the same non-negative infra-marginal price externality $G$ pushes the constrained
optimum \emph{below} the private cutoff (relative to $k_D^{\textup{eq}}$, where the
marginal-type term dominates) yet \emph{above} the first-best (relative to
$k_D^{\textup{FB}}$, where the marginal-type term is zero and only $G$ remains).
The earlier draft's slip conflated these two reference points; the directions are
consistent once the FOC is evaluated at each cutoff in turn.
\end{proof}

\paragraph{Scope.}
The disclosure ordering is conditional on the named hypotheses (i)--(ii), of which
the baseline partial probe corroborates only (i)
($k_D^{\star}<k_D^{\textup{eq}}$, interior, at $\kappa\in\{0.2,0.5,0.8\}$). The
liquidity wedge is approximately zero at baseline
($\kappa^{\ast}\approx\kappa^{\dagger}\approx0.59$). Phase-0 shows the relevant
$\Delta^{\min}$ and $W$ profiles flip sign across the parameter map, so neither an
unconditional non-monotone-welfare claim nor an unconditional wedge orientation is
asserted; the single-peakedness used to convert signed gradients into argmax
orderings is a Numerical Regularity carried inside the proposition statements. All
welfare numbers reproduce via
\texttt{quality\_reports/fixes/D2\_welfare\_verify3.py} against
\texttt{numerical.solver.solve\_valid} (gate $10^{-5}$) and
\texttt{numerical.model.compute\_welfare} (numpy $\ge2$ after the
\texttt{np.trapezoid} patch).

% --------------------------- END WELFARE (replacement) -----------------------


% =============================================================================
% REPLACEMENT FIGURE CAPTION  (replace draft_v2.tex line 1428 with this).
% Removes the OLD "$W$ inherits the nonmonotone shape driven by fundamental value
% creation" / "potential tension" wording (the refuted inherent-hump mechanism).
% =============================================================================
%
% \caption{\textbf{Welfare and the liquidity hump.} Expected total surplus $W$
% and its components---blockholder payoff $W_B$, minority gains $\Delta^{\min}$,
% and bidder surplus $W_{\textup{bid}}$---against noise-trading intensity
% $\kappa$. By Lemma~\ref{lem:endpoints} minority value creation is
% endpoint-symmetric, $\Delta^{\min}(0)=\Delta^{\min}(1)>0$: it does \emph{not}
% vanish at the extremes. The interior hump is \emph{conditional}---it obtains
% under condition~(W$^{\ast}$) and inverts to a trough when the condition
% fails---and the total $W$ peaks at $\kappa^{\ast}$, which may differ from the
% minority optimum $\kappa^{\dagger}$ (the planner's wedge of
% Proposition~\ref{prop:planner-wedge}). At the baseline calibration the two
% optima nearly coincide.}

% =============================================================================
% BIB ENTRY  (add to bibliography.bib if LevitMalenkoMaug2024 is absent).
% =============================================================================
%
% @article{LevitMalenkoMaug2024,
%   author  = {Levit, Doron and Malenko, Nadya and Maug, Ernst},
%   title   = {Trading and Shareholder Democracy},
%   journal = {The Journal of Finance},
%   year    = {2024},
%   volume  = {79},
%   number  = {1},
%   pages   = {257--304}
% }
